\documentclass[12pt]{article}

\usepackage[spanish]{babel}
\usepackage[utf8]{inputenc}
\usepackage{amsmath, amssymb}
\usepackage{booktabs}
\usepackage{geometry}
\usepackage{longtable}
\usepackage{hyperref}

\geometry{margin=2.5cm}

\title{Reporte del Modelo Matemático:\\
Problema de Ubicación de Carga de eBuses (CLP y RCLP)}
\author{}
\date{}

\begin{document}

\maketitle

\section*{Introducción}

Este documento presenta la formulación técnica del Problema de Ubicación de Carga (CLP) y su extensión robusta (RCLP). 
El modelo está diseñado para optimizar la infraestructura de carga de una flota de autobuses eléctricos (eBuses) en entornos urbanos, minimizando la inversión necesaria y garantizando la resiliencia operativa ante fallos en las estaciones de carga.

\section{Nomenclatura del Modelo}

La precisión matemática del modelo depende de la correcta definición de sus parámetros y variables.

\subsection{Constantes (Parámetros)}

\begin{longtable}{lll}
\toprule
Símbolo & Descripción & Unidad/Tipo \\
\midrule
$C_{min}$ & Nivel mínimo de energía para prevenir descarga profunda & kWh \\
$C_{max}$ & Capacidad máxima de la batería & kWh \\
$D_{ji}$ & Energía consumida entre paradas $j$ e $i$ & kWh \\
$T_{ji}$ & Tiempo de tránsito entre paradas & Minutos \\
$\alpha$ & Tasa de carga de la estación & kWh/min \\
$\mu$ & Desviación máxima permitida del horario & Minutos \\
$SM$ & Margen de seguridad entre cargas & Minutos \\
$M$ & Constante grande para linealización & Escalar \\
$\beta$ & Tiempo máximo permitido para una carga & Minutos \\
$\psi$ & Tiempo mínimo de conexión & Minutos \\
$\tau_{bi}$ & Hora programada de llegada del bus $b$ a la parada $i$ & Tiempo \\
\bottomrule
\end{longtable}

\subsection{Variables de Decisión}

\begin{longtable}{ll}
\toprule
Símbolo & Descripción \\
\midrule
$c_{bi}$ & Nivel de energía del bus $b$ al llegar a la parada $i$ \\
$e_{bi}$ & Energía añadida en la parada $i$ \\
$ct_{bi}$ & Duración de la sesión de carga \\
$t_{bi}$ & Tiempo real de llegada del bus \\
$\Delta t_{bi}$ & Desviación absoluta respecto al horario $\tau_{bi}$ \\
$x_{bi}$ & Variable binaria: 1 si el bus carga en la parada \\
$x_{st}$ & Variable binaria: 1 si se instala estación en el sitio \\
$Z_{bdij}$ & 1 si buses $b$ y $d$ comparten cargador \\
$z_{bdi}$ & Variable de ordenamiento temporal \\
$Z''_{bdij}$ & Intersección entre rutas primaria y respaldo \\
$z''_{bdi}$ & Ordenamiento en modo robusto \\
\bottomrule
\end{longtable}

\section{Formulación del Problema de Ubicación de Carga (CLP)}

El modelo CLP optimiza la ubicación de cargadores bajo condiciones operativas normales.

\subsection{Restricciones de Capacidad de Batería}

\begin{align}
C_{min} &\leq c_{bi} \quad \forall b \in B, \forall s_i \in S_b \setminus \{s_0\} \tag{1} \\
c_{bi} + e_{bi} &\leq C_{max} \tag{2} \\
c_{bi} &\leq c_{bj} + e_{bj} - D_{ji} \quad j=i-1 \tag{3} \\
\alpha \cdot ct_{bi} &\geq e_{bi} \tag{4}
\end{align}

Estas restricciones controlan los límites físicos de la batería y el balance energético entre paradas.

\subsection{Restricciones de Tiempo y Desviación}

\begin{align}
t_{bi} &\geq t_{bj} + ct_{bj} - T_{ji} \tag{5} \\
t_{bi} &\geq t_{bj} + ct_{bj} + SM \cdot x_{bj} \tag{6} \\
\Delta t_{bi} &\geq t_{bi} - \tau_{bi} \tag{7} \\
\Delta t_{bi} &\geq \tau_{bi} - t_{bi} \tag{8} \\
\Delta t_{bi} &\leq \mu \tag{9} \\
ct_{bi} &\geq \psi \cdot x_{bi} \tag{10}
\end{align}

Estas ecuaciones modelan el impacto de la carga sobre el cronograma de operación.

\subsection{Instalación de Unidades de Carga}

\begin{align}
\beta \cdot x_{bi} &\geq ct_{bi} \tag{11} \\
x_{st_{bi}} &\geq x_{bi} \tag{12}
\end{align}

Estas restricciones vinculan la actividad de carga con la infraestructura instalada.

\subsection{Restricciones de No Solapamiento}

\begin{align}
Z_{bdij} &\leq x_{bi} \tag{13} \\
Z_{bdij} &\leq x_{dj} \tag{14} \\
x_{bi} + x_{dj} &\leq Z_{bdij} + 1 \tag{15} \\
t_{bi} &\geq t_{dj} + ct_{dj} + SM - M \cdot z_{bdi} \tag{16} \\
t_{dj} &\geq t_{bi} + ct_{bi} + SM - M \cdot z_{dbi} \tag{17} \\
z_{bdi} + z_{dbj} - (1 - Z_{bdij}) &\leq 1 \tag{18}
\end{align}

Este conjunto de restricciones evita que dos buses utilicen simultáneamente el mismo cargador.

\subsection{Energía Mínima Requerida y Función Objetivo}

\begin{align}
\sum_{i=0}^{b_n} e_{bi} &\geq \sum_{i=0}^{b_n-1} D_{i,i+1} - C_{max} + C_{min} \tag{19}
\end{align}

Función objetivo:

\begin{equation}
\min \sum_{st \in ST} x_{st}
\tag{20}
\end{equation}

El objetivo es minimizar el número total de estaciones instaladas.

\section{Extensión Robusta del Modelo (RCLP)}

El modelo robusto permite absorber fallos en cargadores mediante rutas alternativas de carga.

\subsection{Instalación de Cargadores de Respaldo}

\begin{equation}
x'_{bj} \leq (1 - x_{bi})
\tag{21}
\end{equation}

Esta restricción evita que una estación sea su propio respaldo.

\subsection{Cambio de Ruta Primaria a Respaldo}

\begin{align}
c_{bi} &\geq c'_{bi} - M(1 - x_{bi}) \tag{22} \\
t_{bi} &\leq t'_{bi} + M(1 - x_{bi}) \tag{23}
\end{align}

Estas ecuaciones garantizan que el bus tenga suficiente energía y tiempo para desviarse hacia un cargador de respaldo.

\subsection{Retorno al Modo Primario}

\begin{align}
c'_{bi} &\geq c_{bi} - M(1 - x'_{bj}) \tag{24} \\
t'_{bi} &\leq t_{bi} + M(1 - x'_{bj}) \tag{25}
\end{align}

Permiten que el bus retome la operación normal tras usar un cargador de respaldo.

\subsection{No Solapamiento entre Rutas Primarias y de Respaldo}

\begin{align}
Z''_{bdij} &\leq x'_{bi} \tag{26} \\
Z''_{bdij} &\leq x_{dj} \tag{27} \\
x'_{bi} + x_{dj} &\leq Z''_{bdij} + 1 \tag{28} \\
t'_{bi} &\geq t_{dj} + ct_{dj} - M z''_{bdi} \tag{29} \\
t_{dj} &\geq t'_{bi} + ct'_{bi} - M z''_{dbj} \tag{30} \\
z''_{bdi} + z''_{dbj} - (1 - Z''_{bdij}) &\leq 1 \tag{31}
\end{align}

Estas restricciones coordinan el uso de cargadores cuando algunos buses operan en modo primario y otros en modo de respaldo.

\end{document}
